% ---- pg-596 （书页 584） ----

\noindent 12.21\quad 将由一个 901 人组成的样本按三个收入类别和四个职业满意度类别进行交叉分类. 所得结果见\hyperref[tab:12.15]{表 12.15}. 试对这组数据完成对应分析, 并解释结果.

\begin{mvtable}
\begin{tabular}{lcccc}
\toprule
& \multicolumn{4}{c}{职业满意度}\\
\cmidrule(lr){2-5}
收入 & 非常不满意 & 有些不满意 & 比较满意 & 非常满意\\
\midrule
$<25\ 000$ 美元 & 42 & 62 & 184 & 207\\
$25\ 000$ 美元$\sim 50\ 000$ 美元 & 13 & 28 & 81 & 113\\
$>50\ 000$ 美元 & 7 & 18 & 54 & 92\\
\bottomrule
\end{tabular}
\phantomsection\label{tab:12.15}\mvtabcap{表 12.15\quad 收入与职业满意度数据}
\end{mvtable}

\noindent 资料来源: Adapted from data in Table 8.2 in Agresti, A., \textit{Categorical Data Analysis} (New York: John Wiley, 1990).

\noindent 12.22\quad 对\hyperref[tab:12.10]{表 12.10} 中的森林数据进行对应分析, 并验证例 12.22 中的\hyperref[fig:12.28]{图 12.28}.

\noindent 12.23\quad 对\hyperref[tab:12.8]{表 12.8} 中的陶器数据作双重信息图. 解释结果. 双重信息图的结果与\hyperref[fig:12.22]{图 12.22} 中的对应分析结果是否一致? 讨论所得答案(用行比例作为考古场地的观测向量).

\noindent 12.24\quad 对\hyperref[tab:12.14]{表 12.14} 中的心理健康和社会经济状况数据构造双重信息图, 并进行解释. 双重信息图的结果与练习 12.20 中的对应分析结果是否一致? 讨论所得答案(用列比例作为在各种状态下的观测向量).

\noindent 12.25\quad 利用\hyperref[tab:12.13]{表 12.13} 中的考古数据确定度量和非度量的二维标度变换图(见练习 12.19). 给定每个图中各点的坐标, 完成普罗克鲁斯特斯分析, 解释所得结果.

\noindent 12.26\quad \hyperref[tab:8.7]{表 8.7} 中是马里家庭农场数据(见练习 8.28). 去掉离群值 25, 34, 69 和 72, 数据集中还剩 $n=72$ 个观测. 以家庭农场间的欧氏距离作为相似性度量, 用单连接和沃德方法对农场作聚类. 构造连接树图并比较结果. 是否可以看到若干相互区别的农场类?

\noindent 12.27\quad 用标准化的数据重复练习 12.26. 未标准化数据与标准化数据所得结果之间是否有差别? 作出解释.

\noindent 12.28\quad 用\hyperref[tab:8.7]{表 8.7} 中的马里家庭农场数据, 并去掉离群值 25, 34, 69 和 72, 用 $K$ 均值算法作聚类, $K=5$ 和 6. 将结果与 12.26 中的结果比较. 问: 5 或 6 是正确的类数吗? 给出讨论.

\noindent 12.29\quad 用标准化数据重复练习 12.28. 未标准化数据与标准化数据所得结果之间是否有差别? 作出解释.

\noindent 12.30\quad 某公司计划作一次邮件市场宣传. 每寄一个邮件的成本是 1 美元. 该公司掌握 10 万名顾客的信息. 用下面提供的信息构建累计提升图并解释.

总体响应率: 假定我们除知道总体响应率为 $20\%$ 外没有其他模型. 也就是说: 如果给 10 万名顾客中的每位都寄一份邮件(成本为 10 万美元)的话, 我们大约能收回 2 万份响应.

响应模型结果: 假定有一个响应模型可预测谁将回应市场宣传. 我们用此模型对 10 万名顾客打分, 并只联系前 10 000, 前 20 000, ……名顾客, 并预测正响应(回应)的数目. 模型预测的结果列于下表.

\begin{mvtable}
\begin{tabular}{rrr|rrr}
\toprule
成本(美元) & 联系顾客数 & 正响应数 & 成本(美元) & 联系顾客数 & 正响应数\\
\midrule
10 000 & 10 000 & 6 000 & 60 000 & 60 000 & 18 000\\
20 000 & 20 000 & 10 000 & 70 000 & 70 000 & 18 800\\
30 000 & 30 000 & 13 000 & 80 000 & 80 000 & 19 400\\
40 000 & 40 000 & 15 800 & 90 000 & 90 000 & 19 800\\
50 000 & 50 000 & 17 000 & 100 000 & 100 000 & 20 000\\
\bottomrule
\end{tabular}
\end{mvtable}
